\subsection{Overview}

\textit{[This section will be completed in future work after experimental implementation. Placeholder structure provided below.]}

\subsection{Solution Quality Comparison}

\textit{[To be completed: Tables and figures showing:]}

\begin{itemize}
    \item Table: Expected NPV by method (RL, Greedy, MIP, SP, Perfect Info)
    \item Figure: NPV distribution boxplots
    \item Table: Budget utilization rates
    \item Figure: Risk-return tradeoff (NPV vs. variance)
\end{itemize}

\textbf{Expected findings:}
\begin{itemize}
    \item RL outperforms deterministic MIP by $\sim$15--20\% in expected NPV
    \item RL matches or exceeds stochastic programming performance
    \item RL achieves $\sim$90--95\% of perfect-information upper bound
\end{itemize}

\subsection{Computational Performance}

\textit{[To be completed: Tables and figures showing:]}

\begin{itemize}
    \item Table: Solve time by portfolio size and method
    \item Figure: Scalability plot (log-log scale)
    \item Table: Inference latency statistics (mean, P95, P99)
\end{itemize}

\textbf{Expected findings:}
\begin{itemize}
    \item RL inference: 100--200ms per decision (constant with $N$)
    \item MIP solve time: 15--45 minutes for $N=10$, exponential growth
    \item Speedup factor: 10,000$\times$ for $N=10$, increasing with $N$
\end{itemize}

\subsection{Sensitivity Analysis}

\textit{[To be completed: Analysis of performance under varying conditions:]}

\textbf{Uncertainty Level:}
\begin{itemize}
    \item RL maintains robust performance across low/medium/high uncertainty
    \item Deterministic MIP degrades significantly under high uncertainty
    \item SP performance depends on scenario tree quality
\end{itemize}

\textbf{Budget Constraint Tightness:}
\begin{itemize}
    \item RL learns effective prioritization under tight budgets
    \item Greedy heuristic fails to optimize across projects
    \item MIP optimal under loose budgets but slow to compute
\end{itemize}

\textbf{Portfolio Composition:}
\begin{itemize}
    \item Performance consistent across 60/40 domestic/international mix
    \item Sensitivity to extreme compositions (80/20, 20/80) to be analyzed
\end{itemize}

\subsection{Ablation Studies}

\textit{[To be completed: Component-wise analysis:]}

\begin{itemize}
    \item Impact of rolling horizon length $H \in \{6, 12, 18, 24\}$
    \item Effect of reward function components (NPV, utilization, penalties)
    \item Importance of SPI-CPI correlation in training data
    \item Value of attention mechanism (if implemented)
\end{itemize}

\subsection{Discussion}

\textit{[To be completed: Interpretation of results, including:]}

\textbf{Why RL Outperforms Classical Methods:}
\begin{itemize}
    \item Learns robust policies from stochastic simulations
    \item Adapts to partial observability without explicit transition probabilities
    \item Generalizes across portfolio configurations
\end{itemize}

\textbf{Computational Advantage:}
\begin{itemize}
    \item One-time training cost amortized over many deployments
    \item Inference time independent of portfolio size (within trained range)
    \item Enables real-time scenario analysis and Monte Carlo simulation
\end{itemize}

\textbf{Limitations:}
\begin{itemize}
    \item Performance bounded by training data distribution
    \item Requires fine-tuning for company-specific patterns
    \item Black-box nature may limit practitioner trust (addressed by interpretability analysis)
\end{itemize}
